% Source: Gerzensee "Optimization" slides 3-31 (Problem 1 and Remark 12.1;
% binding versus slack constraints and why one identifies them first; the
% overview diagram of implications; Problems 2-5, the Lagrangian, the first-order
% conditions, complementary slackness and the saddle-point formulation;
% Proposition 12.1 and Exercise 12.3; Theorem 12.1, the Kuhn-Tucker theorem with
% its concavity and Slater hypotheses; the Cobb-Douglas consumer problem worked
% through, the multiplier computed two ways, and the envelope calculation).
% Supplementary: Fukuda (2026), Ch. 12 (pp. 423-464) and Sec. 13.4.1
% (Prop. 13.9-13.10, envelope theorems, pp. 479-484); Thm 10.9 (Farkas-Minkowski)
% supplies the proof of Theorem 12.1.
% Code: Supplementary/codes/consumer_choice_fmincon.m, consumer_choice_sqp_Octave.m,
% consumer_choice_python.py, two_periods_consumption_*.{m,py},
% Nash_bargaining_*.{m,py}.
% External references actually cited in this chapter: Boyd-Vandenberghe
% (2004), Mangasarian (1994), Milgrom-Segal (2002), Nash (1950).
% Reported program outputs were obtained by executing the listed files.

\chapter{Constrained Optimisation: Lagrange and Kuhn--Tucker}
\label{ch:optimization}

\section{The Programming Problem}\label{sec:opt-formulation}

\begin{definition}[The programming problem]\label{def:program}
	Let $f,g_1,\dots,g_n\colon\R^m\to\R$. The \dfnas{programming problem}{programming
		problem}, henceforth \dfnas{Problem 1}{Problem 1}, is
	\begin{equation}\label{eq:problem-1}
		\max_{x\in\R^m}\ f(x)
		\qquad\st\quad g_j(x)\geq0\ \text{ for each }j\in\{1,\dots,n\}.
	\end{equation}
	The inequalities are the \dfnas{constraint}{constraints}, $f$ is the
	\dfn{objective function}, and
	\[
		D\eqdef\{x\in\R^m : g_j(x)\geq0\ \text{for each }j\}
	\]
	is the \dfnas{domain of a program}{domain} of the program. The program is
	\dfnas{bounded program}{bounded} if $f$ is bounded above on $D$.
\end{definition}

The sign convention $g_j\geq0$ is not restrictive. A constraint $h(x)\leq0$ is
$-h(x)\geq0$, and an equality constraint $h(x)=0$ is the pair
$h(x)\geq0$ and $-h(x)\geq0$; nothing else is needed
\autocite[slide 4]{fukuda2026optimization}. The convention is chosen so that the
multipliers below are non-negative, and it must be applied consistently: writing
a budget constraint as $m-\ip{p}{x}\geq0$ rather than $\ip{p}{x}-m\geq0$ is what
makes the multiplier on wealth positive rather than negative.

\begin{definition}[Binding and slack]\label{def:binding}
	The constraint $g_j(x)\geq0$ is \dfnas{binding constraint}{binding} at
	$x^\ast\in D$ if $g_j(x^\ast)=0$, and \dfnas{slack constraint}{slack} if
	$g_j(x^\ast)>0$.
\end{definition}

\begin{remark}[Identify the binding constraints first]\label{rem:which-bind}
	Before differentiating anything, determine which constraints bind at the
	optimum. The recommendation of \autocite[slides 6--9]{fukuda2026optimization} is
	not a matter of convenience; it is where the economics of a problem is located,
	and three cases recur.

	\emph{Non-negativity.} With $\lim_{c\downarrow0}u'(c)=\infty$ --- an Inada
	condition --- the constraint $c\geq0$ cannot bind at an optimum, since moving to
	$c=\varepsilon$ raises utility without bound relative to the cost. The
	first-order conditions are then equalities and the analysis is standard. When
	instead a non-negativity-type constraint does bind --- a borrowing constraint, a
	short-sale constraint, limited liability, the zero lower bound on nominal rates
	--- the character of the solution changes, and solving the unconstrained
	first-order condition gives the wrong answer rather than an approximation to the
	right one.

	\emph{Incentive constraints.} A mechanism-design or agency problem may carry
	dozens of incentive-compatibility constraints. Establishing which of them bind
	--- typically the local downward constraints under a single-crossing condition
	--- reduces the problem to a tractable one and is the substantive step.

	\emph{Direction of an equality.} If a constraint is known to hold with equality
	at the optimum, it still matters which inequality is imposed. Writing the budget
	set as $\ip{p}{x}\leq m$ rather than $\ip{p}{x}\geq m$ yields a multiplier with
	the sign that makes it the marginal utility of wealth.
\end{remark}

\section{Five Formulations}\label{sec:opt-five}

The Kuhn--Tucker theorem is the statement that five apparently different
problems have the same solutions once the objective and constraints are concave
and a constraint qualification holds. Stating them
separately makes visible exactly which hypothesis is responsible for which
implication \autocite[slides 10--20]{fukuda2026optimization}.

\begin{definition}[Problem 2: local maximum]\label{def:problem-2}
	Find $x\in D$ such that for every $u\in\R^m$ there is $\lambda_0>0$ with
	\[
		0\leq\lambda\leq\lambda_0\ \text{and}\ x+\lambda u\in D
		\quad\Longrightarrow\quad f(x+\lambda u)\leq f(x).
	\]
\end{definition}

The quantifier structure matters: $\lambda_0$ may depend on the direction $u$.
This is weaker than requiring a single ball around $x$ on which $f(x)$ dominates,
and it is the formulation that matches the way local optimality is verified in
economics --- by considering one-dimensional deviations. The one-shot deviation
principle in repeated games and the perturbation of a consumption path in a
growth model are both instances of checking \autoref{def:problem-2} along a
particular family of directions $u$.

\begin{definition}[Problem 3: no improving feasible direction]\label{def:problem-3}
	Suppose $f,g_1,\dots,g_n$ are differentiable. Find $x\in D$ such that
	\[
		\ip{\nabla f(x)}{u}\leq0
		\quad\text{for every }u\in\R^m\text{ satisfying}\quad
		\bigl[\,g_j(x)=0\ \Longrightarrow\ \ip{\nabla g_j(x)}{u}\geq0\,\bigr].
	\]
\end{definition}

In words: every direction that does not immediately violate a binding constraint
fails to increase the objective to first order. The geometry of
\autocite[slides 12--13]{fukuda2026optimization} is that $\nabla f(x)$ makes an
angle of at least $90^{\circ}$ with every such $u$, while each binding
$\nabla g_j(x)$ makes an angle of at most $90^{\circ}$. With a single binding
constraint this forces $\nabla f(x)$ and $\nabla g_1(x)$ to be antiparallel, that
is $\nabla f(x)+y_1\nabla g_1(x)=0$ for some $y_1\geq0$, which is the Lagrange
condition.

\begin{definition}[Problem 4: Lagrange's problem]\label{def:problem-4}
	Suppose $f,g_1,\dots,g_n$ are differentiable. Find $x\in D$ and
	$y=(y_1,\dots,y_n)\in\R^n_+$ such that
	\begin{align}
		\frac{\partial f}{\partial x_i}(x)
		+\sum_{j=1}^{n}y_j\frac{\partial g_j}{\partial x_i}(x) & =0
		\qquad\text{for each }i\in\{1,\dots,m\},\label{eq:kkt-foc}                     \\
		\sum_{j=1}^{n}y_jg_j(x)                                & =0 .\label{eq:kkt-cs}
	\end{align}
\end{definition}

\begin{definition}[Problem 5: saddle point]\label{def:problem-5}
	Define the \dfn{Lagrangian} $F\colon\R^m\times\R^n_+\to\R$ by
	\begin{equation}\label{eq:lagrangian}
		F(x,y)\eqdef f(x)+\sum_{j=1}^{n}y_jg_j(x).
	\end{equation}
	Find $(\bar x,\bar y)\in\R^m\times\R^n_+$ that is a \dfnas{saddle point}{saddle
		point} of $F$:
	\begin{equation}\label{eq:saddle}
		F(x,\bar y)\leq F(\bar x,\bar y)\leq F(\bar x,y)
		\qquad\text{for all }x\in\R^m\text{ and }y\in\R^n_+ .
	\end{equation}
	\nidx{Lagrangian}{$\Lag$, $F(x,y)$, Lagrangian}
\end{definition}

Condition \eqref{eq:saddle} has two halves: $\bar x$ maximises $F(\cdot,\bar y)$
over the \emph{whole} of $\R^m$, with no constraint, and $\bar y$ minimises
$F(\bar x,\cdot)$ over $\R^n_+$. This is the saddle-point interpretation of
\textcite[Sec.~5.4, p.~237]{boyd2004convex}, whose Section~5.5 (p.~241) states
the optimality conditions \eqref{eq:kkt-foc}--\eqref{eq:kkt-cs} in the same
form and whose Section~5.6 (p.~249) develops the sensitivity reading of the
multiplier recovered below in \autoref{thm:envelope}. The first is what makes the Lagrangian method a
method: an unconstrained problem replaces a constrained one. The second encodes
feasibility, as the next proposition shows.

\begin{proposition}[Complementary slackness at a saddle point]\label{prop:cs}
	Let $(\bar x,\bar y)$ solve Problem 5. Then $\bar x\in D$ and
	$\bar y_jg_j(\bar x)=0$ for every $j$; in particular \eqref{eq:kkt-cs} holds.
\end{proposition}

\begin{proof}
	\emph{Feasibility.} Suppose $g_{j_0}(\bar x)<0$ for some $j_0$. Taking
	$y=\bar y+te_{j_0}$ with $t>0$ gives
	$F(\bar x,y)=F(\bar x,\bar y)+tg_{j_0}(\bar x)\to-\infty$ as $t\to\infty$,
	contradicting the right-hand inequality of \eqref{eq:saddle}. So
	$\bar x\in D$.

	\emph{Complementary slackness.} Taking $y=0\in\R^n_+$ in \eqref{eq:saddle} gives
	$F(\bar x,\bar y)\leq F(\bar x,0)=f(\bar x)$, that is
	$\sum_j\bar y_jg_j(\bar x)\leq0$. Every summand is a product of the
	non-negative numbers $\bar y_j$ and $g_j(\bar x)$, so every summand is
	non-negative; a non-negative sum that is at most zero has all terms zero.
\end{proof}

This is Exercise~12.3 of \textcite{fukuda2026notes}
\autocite[slide 19]{fukuda2026optimization}. The economic reading of
$\bar y_jg_j(\bar x)=0$ is exactly the dichotomy on
\autocite[slide 16]{fukuda2026optimization}: either the constraint binds, or its
marginal valuation is zero. A resource in excess supply has price zero, and a
resource with a positive price is exhausted.

\begin{proposition}[The unconditional implications]\label{prop:p5-p1-p2}
	\begin{enumerate}[(i)]
		\item If $(\bar x,\bar y)$ solves Problem 5, then $\bar x$ solves Problem 1.
		\item If $\bar x$ solves Problem 1, then $\bar x$ solves Problem 2.
	\end{enumerate}
	Neither implication requires any hypothesis on $f$ or the $g_j$.
\end{proposition}

\begin{proof}
	(i) By \autoref{prop:cs}, $\bar x\in D$ and $F(\bar x,\bar y)=f(\bar x)$. For
	any $x\in D$, the left inequality of \eqref{eq:saddle} gives
	\[
		f(x)\leq f(x)+\sum_j\bar y_jg_j(x)=F(x,\bar y)\leq F(\bar x,\bar y)=f(\bar x),
	\]
	the first step because $\bar y_j\geq0$ and $g_j(x)\geq0$ for feasible $x$.

	(ii) A global maximiser over $D$ satisfies the local condition of
	\autoref{def:problem-2} with $\lambda_0$ arbitrary.
\end{proof}

This is Proposition~12.1 of \textcite{fukuda2026notes}
\autocite[slide 19]{fukuda2026optimization}. Part (i) is the part used in
practice: exhibiting a saddle point \emph{proves} global optimality with no
concavity, no differentiability and no constraint qualification. The converse
implications are where hypotheses are needed.

\section{The Kuhn--Tucker Theorem}\label{sec:opt-kkt}

\begin{assumption}[Slater's condition]\label{ass:slater}
	There exists $\hat x\in\R^m$ such that $g_j(\hat x)\geq0$ for every $j$, and
	$g_j(\hat x)>0$ for every $j$ for which $g_j$ is not affine.
\end{assumption}

\begin{theorem}[Kuhn--Tucker]\label{thm:kuhn-tucker}
	Suppose $f,g_1,\dots,g_n$ are concave and \autoref{ass:slater} holds.
	\begin{enumerate}[(i)]
		\item Problems 1, 2 and 5 are equivalent: $\bar x$ solves one if and only if it
		      solves the others, where for Problem 5 this means that some
		      $\bar y\in\R^n_+$ makes $(\bar x,\bar y)$ a saddle point.
		\item If in addition $f,g_1,\dots,g_n$ are differentiable, all five problems
		      are equivalent.
	\end{enumerate}
\end{theorem}

\begin{proof}
	By \autoref{prop:p5-p1-p2} it remains to prove that a solution of Problem 2 is a
	solution of Problem 5, and, under differentiability, to relate Problems 3 and 4
	to the rest.

	\emph{Problem 2 $\Rightarrow$ Problem 1 under concavity.} Let $\bar x$ solve
	Problem 2 and let $x\in D$. Concavity of each $g_j$ makes $D$ convex
	(it is an intersection of upper contour sets of concave, hence quasiconcave,
	functions), so $\bar x+\lambda(x-\bar x)\in D$ for $\lambda\in[0,1]$. Applying
	\autoref{def:problem-2} with $u=x-\bar x$ gives, for small $\lambda>0$,
	\[
		f(\bar x)\geq f\bigl(\bar x+\lambda(x-\bar x)\bigr)
		\geq(1-\lambda)f(\bar x)+\lambda f(x),
	\]
	the second inequality by concavity of $f$. Rearranging,
	$\lambda f(\bar x)\geq\lambda f(x)$, so $f(\bar x)\geq f(x)$.

	\emph{Problem 1 $\Rightarrow$ Problem 5.} Let $\bar x$ solve Problem 1 and set
	$\tilde f(x)\eqdef f(x)-f(\bar x)$, which is concave. Optimality of $\bar x$
	says the system
	\[
		\tilde f(x)>0\quad\text{and}\quad g_j(x)\geq0\ (j=1,\dots,n)
	\]
	has no solution. \autoref{ass:slater} is the solvability hypothesis
	\eqref{eq:slater} of the Farkas--Minkowski theorem, so
	\autoref{thm:farkas} supplies $\bar y\in\R^n_+$ with
	\[
		\sup_{x\in\R^m}\Bigl\{\tilde f(x)+\sum_j\bar y_jg_j(x)\Bigr\}\leq0,
		\qquad\text{that is}\qquad
		F(x,\bar y)\leq f(\bar x)\ \text{ for all }x\in\R^m .
	\]
	Evaluating at $x=\bar x$ gives $\sum_j\bar y_jg_j(\bar x)\leq0$; every term is
	non-negative, so all vanish and $F(\bar x,\bar y)=f(\bar x)$. Hence
	$F(x,\bar y)\leq F(\bar x,\bar y)$ for all $x$, the left half of
	\eqref{eq:saddle}. For the right half, $y\in\R^n_+$ gives
	$F(\bar x,y)=f(\bar x)+\sum_jy_jg_j(\bar x)\geq f(\bar x)=F(\bar x,\bar y)$,
	since each $g_j(\bar x)\geq0$.

	\emph{Problems 3, 4 and 5 under differentiability.} If $(\bar x,\bar y)$ is a
	saddle point, then $\bar x$ maximises the differentiable function
	$F(\cdot,\bar y)$ over the open set $\R^m$, so \autoref{thm:fermat-n} gives
	\eqref{eq:kkt-foc}, and \eqref{eq:kkt-cs} is \autoref{prop:cs}: Problem 5
	implies Problem 4. Conversely, if \eqref{eq:kkt-foc}--\eqref{eq:kkt-cs} hold,
	then $F(\cdot,\bar y)$ is concave --- a non-negative combination of concave
	functions, by \autoref{prop:convex-function-operations}(i) applied to $-F$ ---
	with vanishing gradient at $\bar x$, so \autoref{thm:foc-concave-n} makes
	$\bar x$ a global maximiser of $F(\cdot,\bar y)$, and the right half of
	\eqref{eq:saddle} follows as above: Problem 4 implies Problem 5. That Problem 4
	implies Problem 3 is immediate on taking the inner product of
	\eqref{eq:kkt-foc} with an admissible $u$; the converse uses
	\autoref{thm:farkas} applied to the linearised system at $\bar x$.
\end{proof}

This is Theorem~12.1 of \textcite[p.~436]{fukuda2026notes}
\autocite[slide 20]{fukuda2026optimization}. The proof locates each hypothesis.
Concavity of the $g_j$ makes $D$ convex, which is what turns local into global.
Concavity of $f$ is used in the same step and again in
\autoref{thm:foc-concave-n}. Slater's condition is used once, to invoke
\autoref{thm:farkas}, and it is what guarantees a multiplier \emph{exists};
\autoref{eg:slater-fails} exhibits a program in which it does not.

\begin{remark}[Constraint qualifications]\label{rem:constraint-qualification}
	\autoref{ass:slater} is one of several conditions under which multipliers exist;
	they are collectively the \dfnas{constraint
		qualification}{constraint qualifications}. Three points about them.

	First, they are conditions on the \emph{constraint functions}, not on the
	feasible set. The set $\{x\in\R : -x^2\geq0\}=\{0\}$ and the set
	$\{x\in\R : x\geq0,\ -x\geq0\}=\{0\}$ are identical, but the second
	description satisfies \autoref{ass:slater} --- both constraints are affine ---
	while the first does not, as \autoref{eg:slater-fails} showed. A change of
	description changes whether multipliers exist.

	Second, the affine exemption in \autoref{ass:slater} is not cosmetic. Linear
	programs and problems with only budget and non-negativity constraints satisfy
	the qualification automatically, which is why the constraint qualification is
	almost never mentioned in consumer theory and is essential in problems with
	non-linear constraints such as incentive compatibility.

	Third, when the objective and constraints are not concave the relevant
	qualification is local and differential rather than global. The standard
	condition is that the gradients of the binding constraints at $\bar x$ be
	linearly independent, and under it the conclusion is weaker: the first-order
	conditions \eqref{eq:kkt-foc}--\eqref{eq:kkt-cs} are necessary at a local
	maximum, but not sufficient, and Problem 5 is not equivalent to Problem 1. See
	\textcite[Ch.~7]{mangasarian1994nonlinear}.
\end{remark}

\begin{eg}[Failure of the differential constraint qualification]\label{eg:cq-failure}
	Maximise $f(x_1,x_2)=-x_1$ subject to
	\[
		g_1(x)=-x_2\geq0,\qquad g_2(x)=x_2-(1-x_1)^3\geq0 .
	\]
	The two constraints require $(1-x_1)^3\leq x_2\leq0$, and the left inequality
	forces $x_1\geq1$; conversely every $x_1\geq1$ admits
	$x_2\in[(1-x_1)^3,0]$. The feasible set is therefore
	\[
		D=\bigl\{(x_1,x_2) : x_1\geq1,\ (1-x_1)^3\leq x_2\leq0\bigr\},
	\]
	on which $f=-x_1\leq-1$ with equality only at $x_1=1$, where the interval
	$[(1-x_1)^3,0]$ collapses to $\{0\}$. So $\bar x=(1,0)$ is the unique
	\emph{global} maximiser.

	Both constraints bind at $\bar x$, and
	\[
		\nabla g_1(\bar x)=\begin{pmatrix}0\\-1\end{pmatrix},\qquad
		\nabla g_2(\bar x)=\begin{pmatrix}3(1-x_1)^2\\1\end{pmatrix}
		\bigg|_{x_1=1}=\begin{pmatrix}0\\1\end{pmatrix},
	\]
	which are linearly dependent, so the linear-independence qualification of
	\autoref{rem:constraint-qualification} fails. It fails consequentially: since
	$\nabla f(\bar x)=(-1,0)^{\transp}$ has non-zero first component while both
	constraint gradients have first component $0$, no $y\in\R^2_+$ --- indeed no
	$y\in\R^2$ --- satisfies $\nabla f(\bar x)+y_1\nabla g_1(\bar x)+
		y_2\nabla g_2(\bar x)=0$. The first-order conditions \eqref{eq:kkt-foc} have no
	solution at a point that is globally optimal, and a search over Kuhn--Tucker
	points would miss it.

	Two features of this example deserve separating from each other.
	\autoref{ass:slater} is \emph{not} what fails: $x=(2,-\tfrac12)$ gives
	$g_1=\tfrac12>0$ and $g_2=\tfrac12>0$, so the program is strictly feasible. What
	fails is concavity, since $g_2$ has second derivative $-6(1-x_1)>0$ in $x_1$ for
	$x_1>1$, and with it the hypotheses of \autoref{thm:kuhn-tucker}. Outside the
	concave class the relevant qualification is the local one, and here it is
	violated: the feasible set has an outward cusp at $\bar x$ whose linearisation,
	the set $\{u : \ip{\nabla g_1}{u}\geq0,\ \ip{\nabla g_2}{u}\geq0\}
		=\{u : u_2=0\}$, is strictly larger than the cone of directions along which
	$D$ can actually be entered from $\bar x$, namely $\{u : u_1\geq0,\ u_2=0\}$. It
	is that discrepancy, not the geometry of $D$ on its own, that destroys the
	multiplier rule.
\end{eg}

\subsection{Interpretation of the multipliers}

\begin{theorem}[Envelope theorem]\label{thm:envelope}
	Consider the parametrised program
	\[
		V(b,\alpha)\eqdef\max_{x\in\R^m}\ f(x,\alpha)
		\qquad\st\quad g(x,\alpha)=b,
	\]
	with $f,g$ continuously differentiable, and suppose that on a neighbourhood of
	$(\bar b,\bar\alpha)$ the solution $x^\ast(b,\alpha)$ and the multiplier
	$\lambda^\ast(b,\alpha)$ are unique and continuously differentiable. Then
	\begin{equation}\label{eq:envelope}
		\frac{\partial V}{\partial\alpha}
		=\frac{\partial f}{\partial\alpha}\bigl(x^\ast,\alpha\bigr)
		-\lambda^\ast\,\frac{\partial g}{\partial\alpha}\bigl(x^\ast,\alpha\bigr),
		\qquad
		\frac{\partial V}{\partial b}=\lambda^\ast .
	\end{equation}
\end{theorem}

\begin{proof}
	Write $V(b,\alpha)=f(x^\ast,\alpha)+\lambda^\ast\bigl(b-g(x^\ast,\alpha)\bigr)$,
	which is legitimate because the bracket vanishes at the optimum. Differentiate
	with respect to $\alpha$ using \autoref{thm:chain-rule}:
	\[
		\frac{\partial V}{\partial\alpha}
		=\underbrace{\bigl(f_x-\lambda^\ast g_x\bigr)}_{=0\ \text{by \eqref{eq:kkt-foc}}}
		\cdot\frac{\partial x^\ast}{\partial\alpha}
		+\underbrace{\bigl(b-g(x^\ast,\alpha)\bigr)}_{=0\ \text{by feasibility}}
		\cdot\frac{\partial\lambda^\ast}{\partial\alpha}
		+f_\alpha-\lambda^\ast g_\alpha .
	\]
	The first bracket vanishes by the first-order condition and the second by the
	constraint, leaving the last two terms. Differentiating instead with respect to
	$b$ leaves only the term $\lambda^\ast\cdot\partial b/\partial b=\lambda^\ast$.
\end{proof}

This is Proposition~13.9 and the constrained envelope theorem of
\textcite[pp.~479--481]{fukuda2026notes}, and the computation of
\autocite[slide 28]{fukuda2026optimization}. The content is that the indirect
effects, operating through the response $\partial x^\ast/\partial\alpha$ of the
optimal choice, cancel: at an optimum the marginal value of moving $x$ is zero,
so a marginal change in the parameter has only its direct effect. This is what
makes $\lambda^\ast$ the \emph{shadow price} of the constraint: it is the rate at
which the maximised value rises when the constraint is relaxed, and it is
computable without differentiating the solution.

\begin{remark}[Differentiability of $V$ is an assumption, not a conclusion]
	\label{rem:envelope-caveats}
	\autoref{thm:envelope} presumes that $x^\ast$ and $\lambda^\ast$ are
	differentiable in the parameters, which is itself a conclusion of the implicit
	function theorem (\autoref{thm:implicit-function}) applied to the first-order
	conditions, and requires the bordered Hessian to be non-singular at the optimum.
	When the solution set is not a singleton or when the active set changes, $V$ can
	have a kink and \eqref{eq:envelope} fails as stated. The general form due to
	\textcite{milgrom2002envelope} replaces the derivative by one-sided derivatives
	and requires only that $f_\alpha(x^\ast,\alpha)$ exist at a single maximiser:
	$V'(\alpha^-)\leq f_\alpha(x^\ast,\alpha)\leq V'(\alpha^+)$ whenever the
	one-sided derivatives exist, with equality when $V$ is differentiable. That
	version is what applies to mechanism design, where $X$ may be a discrete set
	with no differentiable structure at all.
\end{remark}

\section{Economic Application: The Consumer's Problem}\label{sec:opt-consumer}

\paragraph{Problem.} With $f(x_1,x_2)=x_1^{\alpha}x_2^{1-\alpha}$,
$\alpha\in(0,1)$, and constraints
\[
	g_1(x)=m-p_1x_1-p_2x_2\geq0,\qquad g_2(x)=x_1\geq0,\qquad g_3(x)=x_2\geq0,
\]
with $p_1,p_2,m>0$, solve \eqref{eq:problem-1}.

\begin{proposition}\label{prop:cobb-douglas-demand}
	The problem has the unique solution
	\begin{equation}\label{eq:cobb-douglas-demand}
		x_1^\ast=\frac{\alpha m}{p_1},\qquad x_2^\ast=\frac{(1-\alpha)m}{p_2},
	\end{equation}
	with multipliers $y_2^\ast=y_3^\ast=0$ and
	\begin{equation}\label{eq:cobb-douglas-multiplier}
		y_1^\ast=\frac{(x_1^\ast)^{\alpha}(x_2^\ast)^{1-\alpha}}{m}
		=\frac{v(p,m)}{m},
	\end{equation}
	where $v(p,m)$ is the indirect utility. Expenditure shares are constant:
	$p_1x_1^\ast/m=\alpha$ and $p_2x_2^\ast/m=1-\alpha$.
\end{proposition}

\begin{proof}
	\emph{Existence and uniqueness.} The budget set is compact
	(\autoref{eg:budget-compact}) and non-empty, and $f$ is continuous, so a
	solution exists by \autoref{cor:weierstrass-continuous}. On $\R^2_{++}$ the
	function $f$ is strictly quasiconcave, so the solution is unique by
	\autoref{prop:strict-quasi-unique}
	\autocite[slide 21]{fukuda2026optimization}.

	\emph{Which constraints bind.} If $x_1^\ast=0$ then $f=0$, while the feasible
	point $(m/(2p_1),m/(2p_2))$ gives $f>0$; so $g_2$ is slack, and symmetrically
	$g_3$. The budget constraint binds: $f$ is strictly increasing in each argument
	on $\R^2_{++}$, so an unspent unit of wealth could be spent to raise utility.

	\emph{First-order conditions.} With $y_2=y_3=0$ by \autoref{prop:cs},
	\eqref{eq:kkt-foc} reads
	\[
		\alpha (x_1^\ast)^{\alpha-1}(x_2^\ast)^{1-\alpha}=y_1p_1,
		\qquad
		(1-\alpha)(x_1^\ast)^{\alpha}(x_2^\ast)^{-\alpha}=y_1p_2 .
	\]
	Dividing eliminates $y_1$ and gives
	$\dfrac{x_2^\ast}{x_1^\ast}=\dfrac{1-\alpha}{\alpha}\cdot\dfrac{p_1}{p_2}$,
	that is $p_2x_2^\ast=\frac{1-\alpha}{\alpha}p_1x_1^\ast$. Substituting into the
	binding budget constraint $p_1x_1^\ast+p_2x_2^\ast=m$ yields
	$p_1x_1^\ast\bigl(1+\frac{1-\alpha}{\alpha}\bigr)=m$, hence
	\eqref{eq:cobb-douglas-demand}.

	\emph{Multiplier.} From the first condition,
	\[
		y_1^\ast=\frac{\alpha(x_1^\ast)^{\alpha-1}(x_2^\ast)^{1-\alpha}}{p_1}
		=\frac{\alpha(x_1^\ast)^{\alpha}(x_2^\ast)^{1-\alpha}}{p_1x_1^\ast}
		=\frac{(x_1^\ast)^{\alpha}(x_2^\ast)^{1-\alpha}}{m},
	\]
	using $p_1x_1^\ast=\alpha m$. The second condition gives the same value, using
	$p_2x_2^\ast=(1-\alpha)m$.

	\emph{Verification via Problem 5.} $f$ is concave on $\R^2_+$ and the $g_j$ are
	affine, so \autoref{ass:slater} holds automatically and
	\autoref{thm:kuhn-tucker} applies; alternatively, the candidate with its
	multipliers satisfies \eqref{eq:kkt-foc}--\eqref{eq:kkt-cs}, which by the last
	part of that proof makes it a saddle point, and
	\autoref{prop:p5-p1-p2}(i) then certifies global optimality without any further
	argument.
\end{proof}

\paragraph{Reading the multiplier.} Formula
\eqref{eq:cobb-douglas-multiplier} says $y_1^\ast=v(p,m)/m$, and since Cobb--Douglas
indirect utility is $v(p,m)=m\cdot(\alpha/p_1)^{\alpha}((1-\alpha)/p_2)^{1-\alpha}$,
the multiplier is independent of $m$: the marginal utility of wealth is constant,
because the utility function is homogeneous of degree one in the bundle and hence
in wealth. \autoref{thm:envelope} confirms this directly,
$\partial v/\partial m=y_1^\ast$, which is the computation of
\autocite[slides 28--29]{fukuda2026optimization} in which every term involving
$\diff x_i^\ast/\diff m$ cancels against a first-order condition.

The constant expenditure shares are the economic signature of Cobb--Douglas
preferences and are the reason the functional form is used: demand for each good
is independent of the price of the other, income elasticities are unity, and the
own-price elasticity is exactly $-1$. Each of these is a strong restriction, and
all three follow from the single fact that the elasticity of substitution is $1$.

\paragraph{Numerical counterpart.} Three of the accompanying programs solve the
same problem by sequential quadratic programming
\autocite[slides 33--36]{fukuda2026optimization}: \path{consumer_choice_fmincon.m},
\path{consumer_choice_sqp_Octave.m} and \path{consumer_choice_python.py}. All
three fix $\alpha=0.7$, $p_1=p_2=1$ and $m=1$. Run as supplied,
\path{consumer_choice_fmincon.m} returns $x^\ast=(0.700,0.300)$ with objective
value $0.54288$ and constraint multiplier $0.54288$;
\path{consumer_choice_python.py} returns the same allocation with objective
$0.54285$ and multiplier $0.54295$, the discrepancy in the fifth decimal being
the solver tolerance. The closed form \eqref{eq:cobb-douglas-demand} gives
$x^\ast=(\alpha m/p_1,(1-\alpha)m/p_2)=(0.7,0.3)$ and
\[
	v(p,m)=\frac{\alpha^{\alpha}(1-\alpha)^{1-\alpha}}{p_1^{\alpha}p_2^{1-\alpha}}\,m
	=0.7^{0.7}\,0.3^{0.3}=0.542884 ,
\]
which is linear in $m$, so $\partial v/\partial m=v(p,m)/m=0.542884$, matching
\eqref{eq:cobb-douglas-multiplier}: the reported multiplier is the shadow price
of income to five decimals, as \autoref{thm:envelope} requires.

The Octave variant, translated into \textsc{Matlab} syntax and run there,
returns $x^\ast=(0.700000,0.300000)$ with utility and multiplier both
$0.542881$, and zero multipliers on the two non-negativity constraints.

The three programs do not share a sign convention, and neither matches
\autoref{def:program}. All three minimise, so each maximisation is passed as
$-u$. \verb|fmincon| takes inequality constraints in the form $c(x)\leq0$ and
returns non-negative multipliers for them; Octave's \verb|sqp| takes them in the
form $h(x)\geq0$, so its $h$ is the negative of \verb|fmincon|'s $c$ and its
multipliers are again non-negative; SciPy's \verb|SLSQP| follows the Octave
convention. The three multipliers therefore agree in sign and magnitude here
only because each program supplies the constraint in the form its own solver
expects. Scaling matters as well: replacing $p_1x_1+p_2x_2\leq m$ by
$2p_1x_1+2p_2x_2\leq2m$ leaves the solution unchanged and halves the reported
multiplier, so a multiplier is a shadow price only relative to the units in
which its constraint is written.

\section{Economic Application: Intertemporal Choice}\label{sec:opt-intertemporal}

\paragraph{Problem.} A consumer with wealth $W>0$ chooses consumption in two
periods, facing gross interest rate $R>0$ and discount factor $\beta\in(0,1)$:
\begin{equation}\label{eq:two-period}
	\max_{c_1,c_2\geq0}\ u(c_1)+\beta u(c_2)
	\qquad\st\quad c_1+\frac{c_2}{R}\leq W,
\end{equation}
with $u$ strictly increasing, strictly concave, differentiable and satisfying
$\lim_{c\downarrow0}u'(c)=\infty$.

\begin{proposition}[Euler equation]\label{prop:euler-two-period}
	Problem \eqref{eq:two-period} has a unique solution, it is interior, the budget
	constraint binds, and the solution satisfies
	\begin{equation}\label{eq:euler-equation}
		u'(c_1^\ast)=\beta R\,u'(c_2^\ast),
	\end{equation}
	with multiplier $\lambda^\ast=u'(c_1^\ast)$. For
	$u(c)=c^{1-\gamma}/(1-\gamma)$ this gives
	$c_2^\ast/c_1^\ast=(\beta R)^{1/\gamma}$ and
	\[
		c_1^\ast=\frac{W}{1+\beta^{1/\gamma}R^{(1-\gamma)/\gamma}} .
	\]
\end{proposition}

\begin{proof}
	The feasible set is compact and convex and the objective continuous and strictly
	concave, so a unique solution exists by
	\autoref{cor:weierstrass-continuous} and \autoref{prop:strict-quasi-unique}. The
	Inada condition rules out $c_i^\ast=0$: at such a point the derivative of the
	objective in $c_i$ is $+\infty$ while the marginal cost is finite, so a small
	increase is feasible and improving. Monotonicity makes the budget constraint
	bind. The constraints are affine, so \autoref{ass:slater} holds and
	\autoref{thm:kuhn-tucker} applies. The Lagrangian
	$F=u(c_1)+\beta u(c_2)+\lambda(W-c_1-c_2/R)$ gives
	$u'(c_1^\ast)=\lambda^\ast$ and $\beta u'(c_2^\ast)=\lambda^\ast/R$;
	eliminating $\lambda^\ast$ yields \eqref{eq:euler-equation}. With CRRA utility,
	$u'(c)=c^{-\gamma}$ and \eqref{eq:euler-equation} reads
	$(c_1^\ast)^{-\gamma}=\beta R(c_2^\ast)^{-\gamma}$, that is
	$c_2^\ast=(\beta R)^{1/\gamma}c_1^\ast$; substituting into the binding budget
	constraint gives the stated $c_1^\ast$.
\end{proof}

\paragraph{Content of \eqref{eq:euler-equation}.} The marginal utility cost of
saving one more unit today equals its discounted marginal utility benefit
tomorrow, scaled by the gross return. The consumption growth rate
$(\beta R)^{1/\gamma}$ is increasing in $R$ with elasticity $1/\gamma$, the
elasticity of intertemporal substitution. The effect of $R$ on $c_1^\ast$ is
ambiguous in sign: differentiating the closed form of
\autoref{prop:euler-two-period}, $\diff c_1^\ast/\diff R>0$ if and only if
$\gamma>1$, the income effect dominating the substitution effect. The multiplier
$\lambda^\ast=u'(c_1^\ast)$ is the marginal utility of wealth, so by
\autoref{thm:envelope} $\partial v/\partial W=u'(c_1^\ast)$, where $v(W)$ denotes
the maximum value of \eqref{eq:two-period}; the same identity in the
infinite-horizon case is the Benveniste--Scheinkman formula of
\autoref{ch:dynamic}.

The corresponding programs are \path{two_periods_consumption_fmincon.m},
\path{two_periods_consumption_sqp_Octave.m} and
\path{two_periods_consumption_python.py}
\autocite[slides 37--40]{fukuda2026optimization}. They take $\beta=0.9$,
$R=1.05$, $W=1$ and logarithmic utility, the case $\gamma=1$ in which the two
offsetting effects cancel and $c_1^\ast=W/(1+\beta)$ is independent of $R$. Run
as supplied, both the \textsc{Matlab} and the Python version return
$c^\ast=(0.526,0.497)$ with objective $-1.27044$ and multiplier $1.90000$. The
closed form gives $c_1^\ast=1/1.9=0.526316$ and
$c_2^\ast=\beta R W/(1+\beta)=0.497368$, lifetime utility
$\log c_1^\ast+\beta\log c_2^\ast=-1.270437$, and
$\lambda^\ast=u'(c_1^\ast)=1/c_1^\ast=1+\beta=1.9$ exactly. The Octave variant,
translated into \textsc{Matlab} syntax and run there, returns
$c^\ast=(0.526316,0.497368)$, lifetime utility $-1.270436$ and multiplier
$1.900000$.

\section{Economic Application: Nash Bargaining}\label{sec:opt-nash}

\paragraph{Problem.} Two players divide a surplus. The set of feasible utility
pairs is $S\subseteq\R^2$, compact and convex, and the disagreement point is
$d\in S$. The \dfnas{Nash bargaining solution}{Nash bargaining solution} solves
\begin{equation}\label{eq:nash-bargaining}
	\max_{(u_1,u_2)\in S}\ (u_1-d_1)(u_2-d_2)
	\qquad\st\quad u_1\geq d_1,\ u_2\geq d_2 .
\end{equation}

\begin{proposition}\label{prop:nash-bargaining}
	Suppose $S$ is compact and convex and there is $u\in S$ with $u\gg d$. Then
	\eqref{eq:nash-bargaining} has a unique solution $u^\ast$, it satisfies
	$u^\ast\gg d$, and if $S=\{u\in\R^2_+ : u_1/a_1+u_2/a_2\leq1\}$ with
	$a_1,a_2>0$ and $d=0$, then $u^\ast=(a_1/2,a_2/2)$: the surplus is split
	equally in normalised units.
\end{proposition}

\begin{proof}
	Existence follows from compactness of the feasible set, which is the
	intersection of $S$ with two closed half-spaces, and continuity of the
	objective (\autoref{cor:weierstrass-continuous}). The objective is not concave,
	but its logarithm $\log(u_1-d_1)+\log(u_2-d_2)$ is strictly concave on the
	region where both factors are positive, and by
	\autoref{prop:concave-implies-quasiconcave} the strictly increasing
	transformation preserves the maximiser; uniqueness is
	\autoref{prop:strict-quasi-unique}. The hypothesis $u\gg d$ makes the maximum
	positive, so $u^\ast\gg d$ and the transformation is legitimate at the optimum.

	For the linear frontier with $d=0$, maximise $\log u_1+\log u_2$ subject to
	$u_1/a_1+u_2/a_2\leq1$. This is \autoref{prop:cobb-douglas-demand} with
	$\alpha=\tfrac12$, $p_i=1/a_i$ and $m=1$, giving $u_i^\ast=a_i/2$.
\end{proof}

Maximising a product subject to a linear constraint is literally the
Cobb--Douglas consumer problem, so the equal split of the Nash solution and the
equal expenditure shares of symmetric Cobb--Douglas preferences are one
computation seen twice.

\begin{remark}[A curved frontier]\label{rem:nash-elliptic}
	Linearity of the frontier plays no role in the split. Take
	$S=\{u\in\R^2_+ : u_1^2/a_1^2+u_2^2/a_2^2\leq1\}$ with $a_1,a_2>0$ and $d=0$.
	The set is the intersection of an ellipse's interior with the positive orthant,
	hence compact and convex, and $(a_1,a_2)/\sqrt2\gg0$ lies in it, so
	\autoref{prop:nash-bargaining} gives a unique solution. Maximising
	$\log u_1+\log u_2$ subject to the constraint, the first-order conditions read
	$1/u_i=2\lambda u_i/a_i^2$, so $u_1^2/a_1^2=u_2^2/a_2^2$; the constraint binds,
	each term equals $\tfrac12$, and
	\[
		u^\ast=\Bigl(\frac{a_1}{\sqrt2},\frac{a_2}{\sqrt2}\Bigr).
	\]
	The normalised shares $u_i^\ast/a_i$ are again equal across players. What
	produces this in both cases is the invariance of the Nash product under the
	rescaling $u_i\mapsto u_i/a_i$, which turns either frontier into a symmetric
	one; it is the scale-invariance axiom of \textcite{nash1950bargaining}, not the
	shape of $S$.
\end{remark}

The numerical counterparts are \path{Nash_bargaining_fmincon.m},
\path{Nash_bargaining_sqp_Octave.m} and \path{Nash_bargaining_python.py}
\autocite[slides 41--44]{fukuda2026optimization}. They solve the curved case of
\autoref{rem:nash-elliptic} with $a_1=a_2=1$: the objective passed to the solver
is $-u_1u_2$ and the non-linear constraint is $u_1^2+u_2^2\leq1$, with lower
bounds $u\geq0$ and starting point $(0.1,0.1)$. Run as supplied, both the
\textsc{Matlab} and the Python version return $u^\ast=(0.7071,0.7071)$, and the
Python version reports multiplier $0.50000$; the Octave variant, translated into
\textsc{Matlab} syntax and run there, returns $(0.707107,0.707107)$ with product
$0.500000$ and multiplier $0.500000$. All agree with
$u^\ast=(1/\sqrt2,1/\sqrt2)=(0.707107,0.707107)$ and, from the first-order
condition $u_2^\ast=2\lambda^\ast u_1^\ast$ for the product objective,
$\lambda^\ast=\tfrac12$.

The starting point is not incidental. The gradient of $(u_1-d_1)(u_2-d_2)$
vanishes at $u=d$, so a solver launched from the disagreement point sees a
stationary point and stops there; $(0.1,0.1)$ is chosen to be interior to the
region where both factors are positive. Solving \eqref{eq:nash-bargaining}
directly is also worse conditioned than solving the logarithmic version, the
product objective being flat along the axes.

\section{Chapter Summary}\label{sec:opt-summary}

A programming problem in the standard form \eqref{eq:problem-1} admits five
formulations: global maximum, local maximum, absence of an improving feasible
direction, the Lagrange conditions with complementary slackness, and a saddle
point of the Lagrangian. Two implications hold with no hypotheses at all: a
saddle point yields a global maximum, and a global maximum is a local maximum.
The converses require concavity of the objective and constraints and a constraint
qualification, and under those hypotheses the Kuhn--Tucker theorem makes all five
equivalent, with differentiability needed only for the two formulations that
mention derivatives. The proof of the essential implication --- from global
optimality to the existence of multipliers --- is an application of the
Farkas--Minkowski theorem, itself a consequence of the separating hyperplane
theorem, so the multipliers are the coefficients of a separating hyperplane.

Complementary slackness says a constraint either binds or has zero multiplier.
The multiplier is the shadow price of the constraint, and the envelope theorem
makes this precise: at an optimum, the derivative of the value function with
respect to a parameter is the partial derivative of the Lagrangian, with all
indirect effects through the optimal choice cancelling against the first-order
conditions. Applied to the Cobb--Douglas consumer, the machinery yields constant
expenditure shares and a multiplier equal to indirect utility divided by wealth;
applied to two-period consumption, the Euler equation
$u'(c_1)=\beta Ru'(c_2)$ with the multiplier equal to the marginal utility of
current consumption; and applied to bargaining, the Nash solution, which is
formally the same optimisation as symmetric Cobb--Douglas demand.

\begin{exercise}\label{ex:opt-1}
	A firm maximises $pF(K,L)-rK-wL$ subject to $K\leq\bar K$, with $F$ concave and
	increasing. Write the Kuhn--Tucker conditions, and show that the multiplier on
	the capacity constraint equals $pF_K-r$ when the constraint binds and $0$
	otherwise. Interpret it as the shadow rental rate, and use
	\autoref{thm:envelope} to compute the value of an extra unit of capacity.
\end{exercise}

\begin{exercise}\label{ex:opt-2}
	Consider $\max x_1$ subject to $g_1(x)=x_2-x_1^3\geq0$ and
	$g_2(x)=-x_2\geq0$. Show that $\bar x=(0,0)$ is the unique feasible point with
	$x_1\geq0$ maximal, that both constraints bind there, and that no multipliers
	satisfy \eqref{eq:kkt-foc}. Identify which hypothesis of
	\autoref{thm:kuhn-tucker} fails, and verify that adding the affine constraint
	$x_1\leq0$ makes the program satisfy \autoref{ass:slater}.
\end{exercise}

\begin{exercise}\label{ex:opt-3}
	Derive the expenditure-minimisation dual of
	\autoref{prop:cobb-douglas-demand}: minimise $p_1x_1+p_2x_2$ subject to
	$x_1^{\alpha}x_2^{1-\alpha}\geq\bar u$. Show that the multiplier on the utility
	constraint is the reciprocal of $y_1^\ast$ from
	\eqref{eq:cobb-douglas-multiplier} evaluated at the corresponding wealth, and
	that the expenditure function $e(p,\bar u)$ satisfies
	$\partial e/\partial p_i=x_i^h$, the Hicksian demand. State which part of
	\autoref{thm:envelope} delivers the last identity.
\end{exercise}
